% Supplementary material of the manuscript. Every number comes from generated/, which the
% repository's generation script writes from the committed records; do not type numbers by
% hand. Build: latexmk -pdf supplementary.tex (from this directory).
\documentclass[11pt,a4paper]{article}

\usepackage[british]{babel}
\usepackage[margin=2.2cm]{geometry}
\usepackage{amsmath,amssymb}
\usepackage{graphicx}
\usepackage{booktabs}
\usepackage{array}
\usepackage{longtable}
\usepackage{pdflscape}
\usepackage{flafter}
\usepackage{microtype}
\usepackage[hidelinks]{hyperref}
\renewcommand{\thetable}{S\arabic{table}}
\newcommand{\cstar}{c^{\ast}}

% generated from the committed records; do not edit
\newcommand{\numAmpFineCstar}{1.23}
\newcommand{\numAmpFineDisp}{0.258}
\newcommand{\numAmpFineDispC}{0.0832}
\newcommand{\numAmpFineGapMedian}{0.0911}
\newcommand{\numAmpFineGapMedianC}{0.0418}
\newcommand{\numAmpFineRedMax}{69\%}
\newcommand{\numAmpFineRedMin}{66\%}
\newcommand{\numAmpIdentity}{\ensuremath{2 \times 10^{-12}}}
\newcommand{\numAmpInCstar}{1.00}
\newcommand{\numAmpInDisp}{0.0300}
\newcommand{\numAmpInDispC}{0.0142}
\newcommand{\numAmpInGap}{0.00912}
\newcommand{\numAmpInGapC}{0.00740}
\newcommand{\numAmpRatioBoth}{5.9}
\newcommand{\numAmpRatioFineOnly}{2.8}
\newcommand{\numAmpRows}{4,608}
\newcommand{\numAuditItems}{22 of 22}
\newcommand{\numChebThree}{0.055}
\newcommand{\numChebTwo}{0.043}
\newcommand{\numCmAdvMax}{99.8\%}
\newcommand{\numCmAdvMin}{71.3\%}
\newcommand{\numCmDispSpread}{1.27}
\newcommand{\numCmFreeDisp}{0.0634}
\newcommand{\numCmFreeDispExcessAbs}{25.6\%}
\newcommand{\numCmFreeDispExcessMax}{37\%}
\newcommand{\numCmFreeDispExcessMin}{15\%}
\newcommand{\numCmFreeDispFour}{0.0634}
\newcommand{\numCmFreeGap}{0.0383}
\newcommand{\numCmFreeGapMedian}{0.0251}
\newcommand{\numCmFreeGapSeeds}{0.0374, 0.0405 and 0.0371}
\newcommand{\numCmFreePeak}{0.128}
\newcommand{\numCmFreeRecall}{0.753}
\newcommand{\numCmFreeVm}{0.123}
\newcommand{\numCmFreeWorse}{0}
\newcommand{\numCmKnormDisp}{0.0547}
\newcommand{\numCmKnormGap}{0.0290}
\newcommand{\numCmKnormGapMedian}{0.0172}
\newcommand{\numCmKnormGapSeeds}{0.0370, 0.0146 and 0.0355}
\newcommand{\numCmKnormPeak}{0.0947}
\newcommand{\numCmKnormRecall}{0.804}
\newcommand{\numCmKnormVm}{0.106}
\newcommand{\numCmKnormWorse}{0}
\newcommand{\numCmKnormWorseSixteen}{11}
\newcommand{\numCmKnormWorseSixtyFour}{1}
\newcommand{\numCmKnormWorseTwoFiveSix}{0}
\newcommand{\numCmLabDisp}{0.0505}
\newcommand{\numCmLabDispFour}{0.0505}
\newcommand{\numCmLabGap}{0.133}
\newcommand{\numCmLabGapMedian}{0.117}
\newcommand{\numCmLabGapSixteen}{20.8}
\newcommand{\numCmLabOverFreeGap}{3.5}
\newcommand{\numCmLabOverFreeGapRoot}{1.9}
\newcommand{\numCmLabOverFreeVm}{1.9}
\newcommand{\numCmLabOverKnormGap}{4.60}
\newcommand{\numCmLabPeak}{0.141}
\newcommand{\numCmLabRecall}{0.668}
\newcommand{\numCmLabVm}{0.235}
\newcommand{\numCmLabWorse}{0}
\newcommand{\numCmLabWorseSixteen}{768}
\newcommand{\numCmLabWorseSixtyFour}{532}
\newcommand{\numCmLabWorseTwoFiveSix}{53}
\newcommand{\numCmMgnBelowLabDisp}{1\%}
\newcommand{\numCmMgnDisp}{0.0499}
\newcommand{\numCmMgnGap}{0.825}
\newcommand{\numCmMgnGapMedian}{0.521}
\newcommand{\numCmMgnOverFreeVm}{4.7}
\newcommand{\numCmMgnOverLabGap}{6.2}
\newcommand{\numCmMgnPeak}{0.321}
\newcommand{\numCmMgnRecall}{0.517}
\newcommand{\numCmMgnVm}{0.578}
\newcommand{\numCmMgnWorse}{187}
\newcommand{\numCmMgnWorseSixtyFour}{766}
\newcommand{\numCmPairs}{768}
\newcommand{\numCmSecondary}{80}
\newcommand{\numCmSolves}{5,120}
\newcommand{\numCmSolvesTrain}{4,096}
\newcommand{\numCmSolvesVal}{1,024}
\newcommand{\numCmVal}{256}
\newcommand{\numCoarsenHigh}{2.5}
\newcommand{\numCoarsenHighCoarse}{0.174}
\newcommand{\numCoarsenHighFine}{0.198}
\newcommand{\numCoarsenHighReduction}{$-$23.7\%}
\newcommand{\numCoarsenHighWiden}{23.7\%}
\newcommand{\numCoarsenLow}{1.8}
\newcommand{\numCoarsenLowCoarse}{0.182}
\newcommand{\numCoarsenLowFine}{0.189}
\newcommand{\numCoarsenLowReduction}{$-$17.6\%}
\newcommand{\numCoarsenLowWiden}{17.6\%}
\newcommand{\numCondThree}{0.012}
\newcommand{\numCondTwo}{0.022}
\newcommand{\numCostAsmFine}{8.07 s}
\newcommand{\numCostAsmOverRitzMax}{27}
\newcommand{\numCostAsmOverRitzMin}{6}
\newcommand{\numCostAsmVal}{900 ms}
\newcommand{\numCostBreakEvenDirect}{84,000}
\newcommand{\numCostBreakEvenMatch}{400,000}
\newcommand{\numCostCgTol}{\ensuremath{10^{-10}}}
\newcommand{\numCostFineDirectSolved}{seven}
\newcommand{\numCostFineDirectTotal}{eight}
\newcommand{\numCostLabelMin}{21}
\newcommand{\numCostRitzFineMax}{1.39 s}
\newcommand{\numCostRitzFineMin}{0.54 s}
\newcommand{\numCostRitzValMax}{144 ms}
\newcommand{\numCostRitzValMin}{34 ms}
\newcommand{\numCostSurrFine}{2.15 s}
\newcommand{\numCostSurrVal}{49 ms}
\newcommand{\numCostThreads}{25}
\newcommand{\numCostThreeFineDirectRatio}{64}
\newcommand{\numCostThreeFineGap}{0.12}
\newcommand{\numCostThreeFineMatchRatio}{1.9}
\newcommand{\numCostThreeValDirectRatio}{27}
\newcommand{\numCostThreeValGap}{0.0085}
\newcommand{\numCostThreeValMatchRatio}{5.2}
\newcommand{\numCostTwoDirectRatio}{1.2}
\newcommand{\numCostTwoGap}{0.029}
\newcommand{\numCostTwoMatchRatio}{2.9}
\newcommand{\numCostWarmFewer}{0.3\%}
\newcommand{\numCostWarmMore}{2.0\%}
\newcommand{\numCrossThree}{\ensuremath{7.43 \times 10^{-3}}}
\newcommand{\numCrossTwo}{\ensuremath{8.39 \times 10^{-3}}}
\newcommand{\numDfourteenAlpha}{\ensuremath{3.6 \times 10^{-4}}}
\newcommand{\numDfourteenDisp}{0.99964}
\newcommand{\numDfourteenGap}{0.99928}
\newcommand{\numDiagFixedMax}{0.22}
\newcommand{\numDiagFixedMin}{0.20}
\newcommand{\numDiagFixedSeeds}{six}
\newcommand{\numDiagScaledMax}{0.52}
\newcommand{\numDiagScaledMin}{0.22}
\newcommand{\numDiagSeeds}{six}
\newcommand{\numEOneBaseDisp}{0.0634}
\newcommand{\numEOneBaseGap}{0.0383}
\newcommand{\numEOneDeltaMax}{0.014}
\newcommand{\numEOneDeltaMin}{0.002}
\newcommand{\numEOneFloor}{0.02}
\newcommand{\numEtwoFineMax}{2.2}
\newcommand{\numEtwoFineMin}{1.8}
\newcommand{\numEtwoGapMax}{3.8}
\newcommand{\numEtwoGapMin}{2.8}
\newcommand{\numEtwoNodesFine}{41,367}
\newcommand{\numEtwoNodesIn}{12,318}
\newcommand{\numEtwoStepMax}{51}
\newcommand{\numEtwoStepMin}{38}
\newcommand{\numEtwoTransformerFineStep}{14.8}
\newcommand{\numFewFactorMax}{5.7}
\newcommand{\numFewFactorMin}{2.3}
\newcommand{\numFewSixteenDisp}{0.0531}
\newcommand{\numFewSixteenGap}{0.0702}
\newcommand{\numFewSixteenPeak}{0.389}
\newcommand{\numFewSixteenScratchDisp}{0.799}
\newcommand{\numFewSixteenVm}{0.166}
\newcommand{\numFewSixtyFourDisp}{0.0298}
\newcommand{\numFewSixtyFourGap}{0.0520}
\newcommand{\numFewSixtyFourPeak}{0.272}
\newcommand{\numFewSixtyFourScratchDisp}{0.193}
\newcommand{\numFewSixtyFourVm}{0.126}
\newcommand{\numFieldBoundMax}{0.85}
\newcommand{\numFieldLoadWorseFree}{1}
\newcommand{\numFieldLoadWorseLab}{126}
\newcommand{\numFieldLoadWorseMgn}{106}
\newcommand{\numFieldLoads}{3,072}
\newcommand{\numFieldRqAboveLab}{89.0\%}
\newcommand{\numFieldRqAboveMgn}{99.9\%}
\newcommand{\numFieldRqFree}{11.9}
\newcommand{\numFieldRqLab}{132}
\newcommand{\numFieldRqMgn}{1,026}
\newcommand{\numFieldVmElemFree}{0.0508}
\newcommand{\numFieldVmElemLab}{0.257}
\newcommand{\numFieldVmElemMgn}{0.226}
\newcommand{\numFieldVmVolFree}{0.0474}
\newcommand{\numFieldVmVolLab}{0.247}
\newcommand{\numFieldVmVolMgn}{0.209}
\newcommand{\numFigFineLoad}{downward body force}
\newcommand{\numFigMedianLoad}{downward body force}
\newcommand{\numFigTwoDDispMax}{0.033}
\newcommand{\numFigTwoDDispMin}{0.026}
\newcommand{\numFigTwoDGapMax}{0.66}
\newcommand{\numFigTwoDGapMin}{0.031}
\newcommand{\numFigTwoDLoad}{shear traction on the top edge}
\newcommand{\numFigWorstFreeC}{0.78}
\newcommand{\numFigWorstFreeDisp}{0.22}
\newcommand{\numFigWorstFreeDispC}{0.045}
\newcommand{\numFigWorstLoad}{downward body force}
\newcommand{\numFigWorstTwiceFree}{12\%}
\newcommand{\numFigWorstTwiceLab}{67\%}
\newcommand{\numFigWorstTwiceMgn}{52\%}
\newcommand{\numFineFreeDisp}{0.258}
\newcommand{\numFineFreeGap}{0.306}
\newcommand{\numFineFreeVm}{0.305}
\newcommand{\numFineKill}{1.5}
\newcommand{\numFineKnnDisp}{0.401}
\newcommand{\numFineLabDisp}{0.257}
\newcommand{\numFineLabGap}{0.985}
\newcommand{\numFineLabVm}{0.537}
\newcommand{\numFineMgnDisp}{0.459}
\newcommand{\numFineMgnGap}{16.3}
\newcommand{\numFineMgnVm}{2.40}
\newcommand{\numFineRatio}{8.6}
\newcommand{\numFineWin}{1.25}
\newcommand{\numGOneFtDisp}{2.0\%}
\newcommand{\numGOneFtGap}{92.6\%}
\newcommand{\numGOneGapReduction}{78.7\%}
\newcommand{\numGOneRetired}{5.8\%}
\newcommand{\numGOneSanityMax}{6.1}
\newcommand{\numGOneSanityMin}{4.5}
\newcommand{\numGOneSanityTableSixteen}{2.4}
\newcommand{\numGOneVmReduction}{65.4\%}
\newcommand{\numGTwoAdvMin}{98.9\%}
\newcommand{\numGTwoFloor}{64}
\newcommand{\numGTwoParity}{$-$36.3\%}
\newcommand{\numGTwoSanityMax}{21.3}
\newcommand{\numGTwoSanitySixteen}{2.71}
\newcommand{\numGTwoSanitySixtyFour}{3.9}
\newcommand{\numGTwoSanityX}{3.0}
\newcommand{\numHOneRel}{71.3\%}
\newcommand{\numHOneTau}{12.4\%}
\newcommand{\numHThreeExRel}{31.9\%}
\newcommand{\numHThreeExTau}{50.2\%}
\newcommand{\numHThreeRel}{24.2\%}
\newcommand{\numHThreeResHigh}{28.0\%}
\newcommand{\numHThreeResLow}{21.9\%}
\newcommand{\numHThreeSeedRel}{1\%, 64\% and 4\%}
\newcommand{\numHThreeSeedsZeroTwo}{2.7\%}
\newcommand{\numHThreeTau}{38.1\%}
\newcommand{\numHTwoARel}{78.2\%}
\newcommand{\numHTwoARootRel}{53.4\%}
\newcommand{\numHTwoATau}{16.4\%}
\newcommand{\numHTwoBRel}{54.9\%}
\newcommand{\numHTwoBTau}{12.3\%}
\newcommand{\numHoldContentDev}{\ensuremath{2 \times 10^{-4}}}
\newcommand{\numHoldHOneRel}{71.3\%}
\newcommand{\numHoldHOneTau}{11.5\%}
\newcommand{\numHoldHThreeRel}{+99.1\%}
\newcommand{\numHoldHThreeTau}{160.2\%}
\newcommand{\numHoldHTwoABaseSeeds}{0.142, 0.147 and 0.121}
\newcommand{\numHoldHTwoAMedHold}{0.0223}
\newcommand{\numHoldHTwoAMedRel}{81.4\%}
\newcommand{\numHoldHTwoAMedTau}{13.4\%}
\newcommand{\numHoldHTwoAMedVal}{0.0210}
\newcommand{\numHoldHTwoANewSeeds}{0.130, 0.0213 and 0.0834}
\newcommand{\numHoldHTwoARatio}{2.69}
\newcommand{\numHoldHTwoARatioHigh}{4.64}
\newcommand{\numHoldHTwoARatioLow}{1.12}
\newcommand{\numHoldHTwoARel}{42.9\%}
\newcommand{\numHoldHTwoASeedZero}{9\%}
\newcommand{\numHoldHTwoATau}{47.4\%}
\newcommand{\numHoldHTwoATop}{254}
\newcommand{\numHoldHTwoATopSeeds}{254.3, 21.6 and 190.4}
\newcommand{\numHoldHTwoBRel}{52.8\%}
\newcommand{\numHoldHTwoBTau}{12.9\%}
\newcommand{\numHoldStates}{21}
\newcommand{\numHoldTThreeLab}{340}
\newcommand{\numHoldTThreeLoads}{8,232}
\newcommand{\numHoldTThreeMgn}{237}
\newcommand{\numHoldTTwoMax}{95.8\%}
\newcommand{\numHoldTTwoMin}{86.3\%}
\newcommand{\numHoldThreeExcl}{34}
\newcommand{\numHoldThreeN}{686}
\newcommand{\numHoldThreeSolves}{2,744}
\newcommand{\numHoldThreeTail}{720}
\newcommand{\numHoldTwoExcl}{33}
\newcommand{\numHoldTwoFreeGap}{0.0393}
\newcommand{\numHoldTwoKnormGap}{0.0781}
\newcommand{\numHoldTwoLoads}{49,332}
\newcommand{\numHoldTwoN}{4,111}
\newcommand{\numHoldTwoPairs}{12,333}
\newcommand{\numHoldTwoSolves}{16,444}
\newcommand{\numHoldTwoTail}{4,144}
\newcommand{\numHoldTwoWorseInstFree}{none}
\newcommand{\numHoldTwoWorseInstKnorm}{9}
\newcommand{\numHoldTwoWorseInstLab}{16}
\newcommand{\numHoldTwoWorseInstMgn}{2,904}
\newcommand{\numHoldTwoWorseLoadFree}{11}
\newcommand{\numHoldTwoWorseLoadKnorm}{21}
\newcommand{\numHoldTwoWorseLoadLab}{57}
\newcommand{\numHoldTwoWorseLoadMgn}{11,399}
\newcommand{\numHoldVAlpha}{5\%}
\newcommand{\numHoldVQuantities}{35}
\newcommand{\numJulyOverCmDispMax}{3.4}
\newcommand{\numJulyOverCmDispMin}{2.6}
\newcommand{\numKappaTauTwo}{\ensuremath{1.5 \times 10^{-7}}}
\newcommand{\numKappaThreeHigh}{\ensuremath{5.2 \times 10^{5}}}
\newcommand{\numKappaThreeLow}{\ensuremath{2.9 \times 10^{4}}}
\newcommand{\numKappaTwoHigh}{\ensuremath{1.5 \times 10^{5}}}
\newcommand{\numKappaTwoLow}{\ensuremath{4.8 \times 10^{3}}}
\newcommand{\numKdClosureGap}{101.3\%}
\newcommand{\numKdClosureVm}{105.0\%}
\newcommand{\numKdFreeDispRel}{25.0\%}
\newcommand{\numKdFreeGapMedian}{0.0075}
\newcommand{\numKdFreeMedMax}{0.0093}
\newcommand{\numKdFreeMedMin}{0.0058}
\newcommand{\numKdFreePairs}{738}
\newcommand{\numKdFreePeakRel}{51.7\%}
\newcommand{\numKdFreeRel}{41.0\%}
\newcommand{\numKdFreeTau}{27.0\%}
\newcommand{\numKdFreeVmRel}{19.4\%}
\newcommand{\numKdFreeWelchHigh}{+16.7\%}
\newcommand{\numKdFreeWelchLow}{$-$98.7\%}
\newcommand{\numKdGapRel}{98.2\%}
\newcommand{\numKdGapTau}{11.7\%}
\newcommand{\numKdHoldFreeRel}{38.0\%}
\newcommand{\numKdHoldFreeTau}{25.6\%}
\newcommand{\numKdHoldGapRel}{98.0\%}
\newcommand{\numKdHoldGapTau}{12.6\%}
\newcommand{\numKdHoldVmRel}{83.0\%}
\newcommand{\numKdHoldVmTau}{10.0\%}
\newcommand{\numKdKnormGapMedian}{0.0044}
\newcommand{\numKdLabGapMedian}{0.25}
\newcommand{\numKdMeansGapRel}{99.98\%}
\newcommand{\numKdMeansGapTau}{197.9\%}
\newcommand{\numKdMeansVmRel}{88.1\%}
\newcommand{\numKdMeansVmTau}{43.7\%}
\newcommand{\numKdMeansVmWelchHigh}{+5.8\%}
\newcommand{\numKdMeansVmWelchLow}{$-$182.0\%}
\newcommand{\numKdSecondary}{30}
\newcommand{\numKdVmRel}{83.6\%}
\newcommand{\numKdVmTau}{10.0\%}
\newcommand{\numKfiveEval}{64}
\newcommand{\numKfiveFiveGap}{0.0430}
\newcommand{\numKfiveFree}{242.5}
\newcommand{\numKfiveNaive}{261.3}
\newcommand{\numKfiveNaiveSaving}{$-$8.7\%}
\newcommand{\numKfiveSaving}{$-$0.9\%}
\newcommand{\numKfiveStartGap}{0.0722}
\newcommand{\numKfiveTol}{\ensuremath{10^{-6}}}
\newcommand{\numKfiveTwentyGap}{0.0260}
\newcommand{\numKfiveZero}{240.4}
\newcommand{\numKthreeBand}{3\%}
\newcommand{\numKthreeBaseRef}{0.201 $\pm$ 0.033}
\newcommand{\numKthreeBaseSeeds}{0.811, 6.99 and 1.66}
\newcommand{\numKthreeGain}{92.5\%}
\newcommand{\numKthreeOutside}{four}
\newcommand{\numKthreeWorse}{three}
\newcommand{\numMgnDispLowerMax}{92.6\%}
\newcommand{\numMgnDispLowerMin}{84.4\%}
\newcommand{\numMgnGapHigherMin}{99.2\%}
\newcommand{\numMgnLowerDispHigherVm}{685}
\newcommand{\numMgnPeakLowerMax}{29\%}
\newcommand{\numMgnPeakLowerMin}{17\%}
\newcommand{\numMgnVmHigher}{767}
\newcommand{\numModesPhrase}{$2 \times 10^{-15}$ in both dimensions}
\newcommand{\numModesThree}{\ensuremath{2 \times 10^{-15}}}
\newcommand{\numModesTwo}{\ensuremath{2 \times 10^{-15}}}
\newcommand{\numNlCmEFree}{0.158}
\newcommand{\numNlCmEKnorm}{0.139}
\newcommand{\numNlCmELab}{0.344}
\newcommand{\numNlCmLabelE}{1.14}
\newcommand{\numNlCmLabelVm}{1.12}
\newcommand{\numNlCmNormE}{2.5}
\newcommand{\numNlCmNormEMean}{2.5}
\newcommand{\numNlCmNormVm}{2.2}
\newcommand{\numNlCmNormVmMean}{2.2}
\newcommand{\numNlCmVmFree}{0.120}
\newcommand{\numNlCmVmKnorm}{0.107}
\newcommand{\numNlCmVmLab}{0.230}
\newcommand{\numNlThreeEFree}{0.0863}
\newcommand{\numNlThreeEKnorm}{0.0666}
\newcommand{\numNlThreeELab}{0.501}
\newcommand{\numNlThreeLabelE}{1.30}
\newcommand{\numNlThreeLabelVm}{1.24}
\newcommand{\numNlThreeLabelWithin}{1.3}
\newcommand{\numNlThreeNormE}{7.5}
\newcommand{\numNlThreeNormEMean}{10.7}
\newcommand{\numNlThreeNormVm}{6.1}
\newcommand{\numNlThreeNormVmMean}{8.4}
\newcommand{\numNlThreeNormVmWord}{six}
\newcommand{\numNlThreeVmFree}{0.0527}
\newcommand{\numNlThreeVmKnorm}{0.0425}
\newcommand{\numNlThreeVmLab}{0.259}
\newcommand{\numNodesFineMax}{71,749}
\newcommand{\numNodesFineMedian}{36,216}
\newcommand{\numNodesFineMin}{15,175}
\newcommand{\numNodesInMax}{17,692}
\newcommand{\numNodesInMedian}{5,617}
\newcommand{\numNodesInMin}{2,195}
\newcommand{\numPairCmFreeKnormGapTwoFiveSix}{72.4\%}
\newcommand{\numPairCmKnormFreeDisp}{68.1\%}
\newcommand{\numPairCmKnormFreeGap}{87.6\%}
\newcommand{\numPairCmKnormFreePeak}{84.9\%}
\newcommand{\numPairCmKnormFreeRecall}{74.1\%}
\newcommand{\numPairCmKnormFreeVm}{83.3\%}
\newcommand{\numPairCmKnormLabDisp}{49.7\%}
\newcommand{\numPairCmKnormLabGap}{99.9\%}
\newcommand{\numPairCmKnormLabPeak}{84.6\%}
\newcommand{\numPairCmKnormLabRecall}{98.6\%}
\newcommand{\numPairCmKnormLabVm}{99.9\%}
\newcommand{\numPairCmLabDisp}{29.9\%}
\newcommand{\numPairCmLabGap}{99.9\%}
\newcommand{\numPairCmLabPeak}{62.6\%}
\newcommand{\numPairCmLabRecall}{96.7\%}
\newcommand{\numPairCmLabVm}{99.9\%}
\newcommand{\numPairCmMgnLabDisp}{55.7\%}
\newcommand{\numPairCmMgnLabGap}{none}
\newcommand{\numPairCmMgnLabPeak}{41.0\%}
\newcommand{\numPairCmMgnLabRecall}{9.8\%}
\newcommand{\numPairCmMgnLabVm}{none}
\newcommand{\numPairCmOpposite}{69.9\%}
\newcommand{\numPairThreeAncDisp}{81.2\%}
\newcommand{\numPairThreeAncGap}{99.7\%}
\newcommand{\numPairThreeAncPeak}{75.3\%}
\newcommand{\numPairThreeAncRecall}{95.7\%}
\newcommand{\numPairThreeAncVm}{99.7\%}
\newcommand{\numPairThreeLabDisp}{78.9\%}
\newcommand{\numPairThreeLabGap}{99.9\%}
\newcommand{\numPairThreeLabPeak}{87.4\%}
\newcommand{\numPairThreeLabRecall}{all}
\newcommand{\numPairThreeLabVm}{99.9\%}
\newcommand{\numPairTwoAncDisp}{36.2\%}
\newcommand{\numPairTwoAncGap}{91.8\%}
\newcommand{\numPairTwoAncPeak}{44.8\%}
\newcommand{\numPairTwoAncRecall}{76.4\%}
\newcommand{\numPairTwoAncVm}{75.1\%}
\newcommand{\numPairTwoLabDisp}{39.1\%}
\newcommand{\numPairTwoLabGap}{all}
\newcommand{\numPairTwoLabPeak}{63.2\%}
\newcommand{\numPairTwoLabRecall}{98.8\%}
\newcommand{\numPairTwoLabVm}{99.3\%}
\newcommand{\numParamsAux}{0.6}
\newcommand{\numParamsGraph}{3.9}
\newcommand{\numParamsTransformer}{6.7}
\newcommand{\numPilotDisp}{0.309}
\newcommand{\numPilotLimit}{0.90}
\newcommand{\numPremiseThree}{0.462}
\newcommand{\numPremiseTwo}{0.464}
\newcommand{\numRemeshFine}{1.27}
\newcommand{\numRemeshGeometries}{16}
\newcommand{\numRemeshInDev}{2\%}
\newcommand{\numRerunAccDev}{\ensuremath{4 \times 10^{-16}}}
\newcommand{\numRerunCpu}{Xeon Gold 6459C}
\newcommand{\numRerunIters}{two}
\newcommand{\numRerunThreads}{16}
\newcommand{\numRhoThree}{0.827}
\newcommand{\numRhoTwo}{0.730}
\newcommand{\numRitzFineBnPeakMax}{62\%}
\newcommand{\numRitzFineBnPeakMin}{7\%}
\newcommand{\numRitzFineFreePeakMax}{92\%}
\newcommand{\numRitzFineFreePeakMin}{79\%}
\newcommand{\numRitzFineFtPeakMax}{98\%}
\newcommand{\numRitzFineFtPeakMin}{62\%}
\newcommand{\numRitzFineGapAll}{99.6\%}
\newcommand{\numRitzFineGapNine}{0.0070}
\newcommand{\numRitzFineGapSix}{0.019}
\newcommand{\numRitzFineLabPeakMax}{17.6\%}
\newcommand{\numRitzHoldFreeDispNineMax}{100\%}
\newcommand{\numRitzHoldFreeDispNineMin}{99.6\%}
\newcommand{\numRitzHoldFreeDispSixMax}{81\%}
\newcommand{\numRitzHoldFreeDispSixMin}{67\%}
\newcommand{\numRitzHoldFreePeakMin}{99.4\%}
\newcommand{\numRitzHoldKnormGapMin}{88\%}
\newcommand{\numRitzHoldKnormPeakMiss}{two}
\newcommand{\numRitzThreeK}{168, 252, 360 and 495}
\newcommand{\numRitzThreeKMax}{495}
\newcommand{\numRitzTwoFreeDispMax}{97\%}
\newcommand{\numRitzTwoFreeDispMin}{80\%}
\newcommand{\numRitzTwoFreeGapEightMax}{66\%}
\newcommand{\numRitzTwoFreeGapEightMin}{61\%}
\newcommand{\numRitzTwoFreeGapSixMax}{20\%}
\newcommand{\numRitzTwoFreeGapSixMin}{16\%}
\newcommand{\numRitzTwoFreePeakMin}{96.8\%}
\newcommand{\numRitzTwoK}{42 and 72}
\newcommand{\numRitzTwoKEight}{72}
\newcommand{\numRitzTwoKSix}{42}
\newcommand{\numRitzValFreeDispNineMax}{100\%}
\newcommand{\numRitzValFreeDispNineMin}{99.6\%}
\newcommand{\numRitzValFreeDispSixMax}{79\%}
\newcommand{\numRitzValFreeDispSixMin}{67\%}
\newcommand{\numRitzValFreePeakMin}{98.4\%}
\newcommand{\numRitzValGapNine}{0.0063}
\newcommand{\numRitzValGapSix}{0.018}
\newcommand{\numRitzValKnormGapMin}{89.8\%}
\newcommand{\numRitzValNineDispAll}{97.6\%}
\newcommand{\numRitzValPeakNine}{0.23}
\newcommand{\numRitzValPeakRatioMax}{4.9}
\newcommand{\numRitzValPeakRatioMin}{2.4}
\newcommand{\numRitzValPeakSix}{0.32}
\newcommand{\numSpecBoundMax}{0.90}
\newcommand{\numSpecFigIndex}{221}
\newcommand{\numSpecFreeMax}{3,622}
\newcommand{\numSpecFreeMin}{348}
\newcommand{\numSpecIeeeLabKnormGap}{11.9}
\newcommand{\numSpecIeeeLabKnormRq}{9.30}
\newcommand{\numSpecIeeeRqFree}{9.12}
\newcommand{\numSpecIeeeRqKnorm}{9.42}
\newcommand{\numSpecIeeeRqLab}{114}
\newcommand{\numSpecIeeeRqMgn}{558}
\newcommand{\numSpecKnormFreeAbove}{49.1\%}
\newcommand{\numSpecKnormFreeDisp}{0.64}
\newcommand{\numSpecKnormFreeDispSeeds}{0.81, 0.39 and 0.81}
\newcommand{\numSpecKnormFreeGap}{0.66}
\newcommand{\numSpecKnormFreeRq}{1.04}
\newcommand{\numSpecLabKnormAbove}{94.9\%}
\newcommand{\numSpecLabKnormAboveSeedGm}{98.8\%}
\newcommand{\numSpecLabKnormDisp}{1.26}
\newcommand{\numSpecLabKnormDispLower}{50.0\%}
\newcommand{\numSpecLabKnormDispSeeds}{0.96, 1.65 and 1.26}
\newcommand{\numSpecLabKnormGap}{9.85}
\newcommand{\numSpecLabKnormRq}{7.80}
\newcommand{\numSpecLabKnormRqSeeds}{8.26, 9.46 and 6.07}
\newcommand{\numSpecLoadCases}{1,024}
\newcommand{\numSpecLoadWorseFree}{1}
\newcommand{\numSpecLoadWorseKnorm}{0}
\newcommand{\numSpecLoadWorseLab}{1}
\newcommand{\numSpecLoadWorseMgn}{724}
\newcommand{\numSpecLoads}{3,072}
\newcommand{\numSpecMgnLabAbove}{93.4\%}
\newcommand{\numSpecMgnLabDisp}{1.00}
\newcommand{\numSpecMgnLabGap}{4.34}
\newcommand{\numSpecMgnLabRq}{4.33}
\newcommand{\numSpecRqFree}{10.3}
\newcommand{\numSpecRqKnorm}{11.3}
\newcommand{\numSpecRqLab}{115}
\newcommand{\numSpecRqMgn}{553}
\newcommand{\numSpecRqStarOverMin}{1.27}
\newcommand{\numSpecShareEucFree}{7.2\%}
\newcommand{\numSpecShareEucKnorm}{6.6\%}
\newcommand{\numSpecShareEucLab}{24\%}
\newcommand{\numSpecShareEucMgn}{49\%}
\newcommand{\numSpecShareEucStar}{0.2\%}
\newcommand{\numSpecShareKFree}{69\%}
\newcommand{\numSpecShareKKnorm}{73\%}
\newcommand{\numSpecShareKLab}{95\%}
\newcommand{\numSpecShareKMgn}{99\%}
\newcommand{\numSpecShareKStar}{1.5\%}
\newcommand{\numSpecSpectralShare}{90\%}
\newcommand{\numSpecSpectralShareMax}{102\%}
\newcommand{\numSpecSpectralShareMin}{82\%}
\newcommand{\numSpecTfFree}{2.7\%}
\newcommand{\numSpecTfKnorm}{2.9\%}
\newcommand{\numSpecTfLab}{0.7\%}
\newcommand{\numSpecTfMgn}{0.05\%}
\newcommand{\numSpecTfStiffFree}{41\%}
\newcommand{\numSpecTfStiffKnorm}{36\%}
\newcommand{\numSpecTfStiffLoads}{670}
\newcommand{\numSpecTfStiffN}{2,010}
\newcommand{\numSpecVmAreaFree}{0.105}
\newcommand{\numSpecVmAreaKnorm}{0.0896}
\newcommand{\numSpecVmAreaLab}{0.227}
\newcommand{\numSpecVmAreaMgn}{0.462}
\newcommand{\numSpecVmAreaRatio}{2.26}
\newcommand{\numSpecVmElemFree}{0.112}
\newcommand{\numSpecVmElemKnorm}{0.0964}
\newcommand{\numSpecVmElemLab}{0.231}
\newcommand{\numSpecVmElemMgn}{0.474}
\newcommand{\numSpecVmElemRatio}{2.22}
\newcommand{\numSqrtKappaMax}{723}
\newcommand{\numSqrtKappaMin}{69}
\newcommand{\numTauSqKappaTwo}{\ensuremath{6 \times 10^{-18}}}
\newcommand{\numTauSquared}{\ensuremath{10^{-12}}}
\newcommand{\numThreadsCgGainMax}{29\%}
\newcommand{\numThreadsCgGainMin}{16\%}
\newcommand{\numThreadsDirectDev}{1.1\%}
\newcommand{\numThreadsDrift}{2.1\%}
\newcommand{\numThreadsFineN}{2}
\newcommand{\numThreadsHost}{25}
\newcommand{\numThreadsTwoCg}{4.3\%}
\newcommand{\numThreadsTwoDirect}{7.2\%}
\newcommand{\numThreadsTwoN}{8}
\newcommand{\numThreadsValN}{8}
\newcommand{\numThreeAncDisp}{0.0470}
\newcommand{\numThreeAncGap}{0.116}
\newcommand{\numThreeAncGapMedian}{0.108}
\newcommand{\numThreeAncOverFreeGap}{12.7}
\newcommand{\numThreeAncOverFreeGapRoot}{3.6}
\newcommand{\numThreeAncPeak}{0.115}
\newcommand{\numThreeAncRecall}{0.730}
\newcommand{\numThreeAncVm}{0.128}
\newcommand{\numThreeAncWorse}{2}
\newcommand{\numThreeAncWorseTwoFiveSix}{4}
\newcommand{\numThreeEvalSolves}{2,048}
\newcommand{\numThreeFewSolves}{256}
\newcommand{\numThreeFreeDisp}{0.0300}
\newcommand{\numThreeFreeDispAdvantage}{36.3\%}
\newcommand{\numThreeFreeDispFour}{0.0300}
\newcommand{\numThreeFreeGap}{0.00912}
\newcommand{\numThreeFreeGapMedian}{0.00763}
\newcommand{\numThreeFreeLoadWorse}{1}
\newcommand{\numThreeFreeLoads}{3,072}
\newcommand{\numThreeFreeOverMgnDisp}{1.8}
\newcommand{\numThreeFreePeak}{0.0807}
\newcommand{\numThreeFreeRecall}{0.793}
\newcommand{\numThreeFreeVm}{0.0554}
\newcommand{\numThreeFreeWorse}{0}
\newcommand{\numThreeGuardTotal}{ten}
\newcommand{\numThreeGuardWithin}{three}
\newcommand{\numThreeKnormDisp}{0.0227}
\newcommand{\numThreeKnormGap}{0.00504}
\newcommand{\numThreeKnormPeak}{0.0386}
\newcommand{\numThreeKnormRecall}{0.858}
\newcommand{\numThreeKnormVm}{0.0440}
\newcommand{\numThreeKnormWorse}{0}
\newcommand{\numThreeLabDisp}{0.0470}
\newcommand{\numThreeLabDispFour}{0.0470}
\newcommand{\numThreeLabGap}{31.2}
\newcommand{\numThreeLabGapMedian}{0.252}
\newcommand{\numThreeLabGapMs}{31.2 $\pm$ 53.5}
\newcommand{\numThreeLabGapSeeds}{0.297, 93.0, 0.365}
\newcommand{\numThreeLabPeak}{0.784}
\newcommand{\numThreeLabRecall}{0.619}
\newcommand{\numThreeLabVm}{0.370}
\newcommand{\numThreeLabWorse}{18}
\newcommand{\numThreeLabWorseTwoFiveSix}{132}
\newcommand{\numThreeLabWorst}{23,727}
\newcommand{\numThreeLabWorstSeed}{1}
\newcommand{\numThreeMedianRatioMax}{33}
\newcommand{\numThreeMedianRatioMin}{14}
\newcommand{\numThreeMgnDisp}{0.0168}
\newcommand{\numThreeMgnGap}{0.447}
\newcommand{\numThreeMgnGapMedian}{0.208}
\newcommand{\numThreeMgnOverFreePeak}{6.0}
\newcommand{\numThreeMgnOverFreeVm}{5.3}
\newcommand{\numThreeMgnPeak}{0.488}
\newcommand{\numThreeMgnRecall}{0.679}
\newcommand{\numThreeMgnVm}{0.296}
\newcommand{\numThreeMgnWorse}{28}
\newcommand{\numThreePairs}{768}
\newcommand{\numThreeSolves}{6,400}
\newcommand{\numThreeTrainSolves}{4,096}
\newcommand{\numThreeVal}{256}
\newcommand{\numTimeBottleneckSeedMax}{2.5}
\newcommand{\numTimeBottleneckSeedMin}{2.5}
\newcommand{\numTimeBottleneckSpeedup}{11}
\newcommand{\numTimeThreeSeed}{27.3}
\newcommand{\numTimeTwoParallel}{2 h 58 min}
\newcommand{\numTwoAdvMax}{98.7\%}
\newcommand{\numTwoAdvMin}{80.4\%}
\newcommand{\numTwoAncDisp}{0.161}
\newcommand{\numTwoAncGap}{0.166}
\newcommand{\numTwoAncGapMedian}{0.0709}
\newcommand{\numTwoAncPeak}{0.246}
\newcommand{\numTwoAncRecall}{0.708}
\newcommand{\numTwoAncSixteenSeeds}{0.225, 0.458, 0.546}
\newcommand{\numTwoAncVm}{0.257}
\newcommand{\numTwoAncWorse}{14}
\newcommand{\numTwoDispSpread}{1.05}
\newcommand{\numTwoFreeDisp}{0.166}
\newcommand{\numTwoFreeDispExcess}{+3.6\%}
\newcommand{\numTwoFreeDispExcessAbs}{3.6\%}
\newcommand{\numTwoFreeDispFour}{0.1658}
\newcommand{\numTwoFreeGap}{0.0817}
\newcommand{\numTwoFreeGapMedian}{0.0493}
\newcommand{\numTwoFreePeak}{0.228}
\newcommand{\numTwoFreeRecall}{0.739}
\newcommand{\numTwoFreeVm}{0.210}
\newcommand{\numTwoFreeWorse}{3}
\newcommand{\numTwoFtDisp}{0.166}
\newcommand{\numTwoFtGap}{0.236}
\newcommand{\numTwoFtGapMedian}{0.156}
\newcommand{\numTwoFtPeak}{0.301}
\newcommand{\numTwoFtRecall}{0.645}
\newcommand{\numTwoFtVm}{0.336}
\newcommand{\numTwoFtWorse}{21}
\newcommand{\numTwoGapSpread}{31}
\newcommand{\numTwoLabDisp}{0.160}
\newcommand{\numTwoLabDispFour}{0.1600}
\newcommand{\numTwoLabGap}{0.418}
\newcommand{\numTwoLabGapMedian}{0.274}
\newcommand{\numTwoLabGapSixteen}{6.26}
\newcommand{\numTwoLabOverFreeGap}{5.1}
\newcommand{\numTwoLabPeak}{0.334}
\newcommand{\numTwoLabRecall}{0.584}
\newcommand{\numTwoLabVm}{0.433}
\newcommand{\numTwoLabWorse}{52}
\newcommand{\numTwoMgnDisp}{0.167}
\newcommand{\numTwoMgnGap}{2.57}
\newcommand{\numTwoMgnGapMedian}{1.35}
\newcommand{\numTwoMgnPeak}{1.32}
\newcommand{\numTwoMgnRecall}{0.322}
\newcommand{\numTwoMgnVm}{1.10}
\newcommand{\numTwoMgnWorse}{481}
\newcommand{\numTwoPairs}{768}
\newcommand{\numTwoRuns}{nine}
\newcommand{\numTwoRunsDisp}{0.1658 $\pm$ 0.0023}
\newcommand{\numTwoRunsGap}{0.0820 $\pm$ 0.0026}
\newcommand{\numTwoSolves}{7,168}
\newcommand{\numTwoSolvesOffline}{5,120}
\newcommand{\numTwoSolvesRun}{2,048}
\newcommand{\numTwoVal}{256}
\newcommand{\numWarmFractionSeven}{15\%}
\newcommand{\numWarmFractionSeventeen}{6\%}
\newcommand{\numWarmFractionTwelve}{9\%}
\newcommand{\numWithinThree}{\ensuremath{6.76 \times 10^{-2}}}
\newcommand{\numWithinTwo}{\ensuremath{7.56 \times 10^{-2}}}
\newcommand{\numWorstFreeDisp}{0.071}
\newcommand{\numWorstFreeGap}{0.0349}
\newcommand{\numWorstFreeVm}{0.10}
\newcommand{\numWorstIndex}{107}
\newcommand{\numWorstLabDisp}{0.032}
\newcommand{\numWorstLabGap}{2.90}
\newcommand{\numWorstLabVm}{0.73}
\newcommand{\numWorstMgnDisp}{0.033}
\newcommand{\numWorstMgnGap}{27.0}
\newcommand{\numWorstMgnVm}{2.4}
\newcommand{\numWorstSeedLabDisp}{0.55}
\newcommand{\numWorstSeedLabPeak}{361}
\newcommand{\numWorstSeedLabVm}{65}


\title{Supplementary material for\\[2pt] \large Training finite-element surrogates without
solutions: the discrete potential energy as exact objective, metric and failure detector}
\author{Ruifeng Cao \and Xidan Song}
\date{}

\begin{document}
\maketitle

\section*{S1. The Ritz projection paired with the networks}

The paper pairs the polynomial Ritz-Galerkin projection of each instance, computed from the
instance's own stiffness matrix and loads without labels or training, with every network's
prediction of the same instance; this baseline was added after the original analyses and was
not pre-registered, as the paper states. Table~\ref{tab:ritzpairs} gives, for each set,
model and degree $p$ of the ladder, the share of the instances on which the projection was
strictly the more accurate, as the lowest and the highest over the three seeds of a trained
network. The sets are the validation instances of the three- and two-dimensional
comparisons, their held-out instances, and the finer three-dimensional mesh, on which the
networks trained in-band were applied zero-shot. The projection has \numRitzThreeK{}
unknowns at $p=6$, 7, 8 and 9 in three dimensions and \numRitzTwoK{} at $p=6$ and 8 in two.

\begin{landscape}
% generated from the committed records; do not edit
\begingroup\small\setlength{\LTcapwidth}{\linewidth}
\begin{longtable}{@{}llcccccc@{}}
\caption{The Ritz projection paired with the networks and baselines: the share of the instances on which the projection was the more accurate, per set and degree $p$.}\label{tab:ritzpairs}\\
\toprule
Set & Model & $p$ & \shortstack{Displacement \\ error} & \shortstack{Relative \\ energy gap} & \shortstack{von Mises \\ error} & \shortstack{Peak stress \\ error} & \shortstack{Critical-region \\ recall} \\
\midrule
\endfirsthead
\caption[]{The Ritz projection paired with the networks and baselines: the share of the instances on which the projection was the more accurate, per set and degree $p$. (continued)}\\
\toprule
Set & Model & $p$ & \shortstack{Displacement \\ error} & \shortstack{Relative \\ energy gap} & \shortstack{von Mises \\ error} & \shortstack{Peak stress \\ error} & \shortstack{Critical-region \\ recall} \\
\midrule
\endhead
\bottomrule
\endfoot
3D, validation & Label-free & 6 & 67\%--79\% & 0.4\%--6\% & 1\%--8\% & 0\%--1\% & 48\%--71\% \\
 &  & 7 & 91\%--98\% & 3\%--25\% & 14\%--31\% & 0\%--1\% & 81\%--96\% \\
 &  & 8 & 98\%--99\% & 11\%--73\% & 37\%--67\% & 0\%--1\% & 95\%--99.6\% \\
 &  & 9 & 99.6\%--100\% & 38\%--100\% & 69\%--98\% & 0\%--2\% & 99.6\%--100\% \\
 & Supervised, displacement & 6 & 93\%--94\% & 100\% & 100\% & 31\%--37\% & 100\% \\
 &  & 7 & 99.6\%--100\% & 100\% & 100\% & 36\%--42\% & 100\% \\
 &  & 8 & 100\% & 100\% & 100\% & 41\%--46\% & 100\% \\
 &  & 9 & 100\% & 100\% & 100\% & 47\%--51\% & 100\% \\
 & Supervised, displacement + energy term & 6 & 81\%--100\% & 100\% & 95\%--100\% & 0.4\% & 93\%--95\% \\
 &  & 7 & 97\%--100\% & 100\% & 100\% & 0.4\% & 99\%--99.6\% \\
 &  & 8 & 100\% & 100\% & 100\% & 1\%--2\% & 100\% \\
 &  & 9 & 100\% & 100\% & 100\% & 2\% & 100\% \\
 & Supervised, energy norm & 6 & 43\%--57\% & 0\%--1\% & 0\%--2\% & 0\% & 23\%--25\% \\
 &  & 7 & 78\%--85\% & 0\%--1\% & 4\%--7\% & 0\% & 51\%--54\% \\
 &  & 8 & 94\%--96\% & 1\%--4\% & 13\%--19\% & 0\% & 61\%--65\% \\
 &  & 9 & 98\%--99\% & 6\%--10\% & 28\%--43\% & 0\% & 74\%--80\% \\
 & Graph network, displacement & 6 & 16\%--19\% & 100\% & 100\% & 30\%--35\% & 99.6\%--100\% \\
 &  & 7 & 56\%--62\% & 100\% & 100\% & 32\%--39\% & 100\% \\
 &  & 8 & 90\%--93\% & 100\% & 100\% & 35\%--43\% & 100\% \\
 &  & 9 & 99\% & 100\% & 100\% & 38\%--48\% & 100\% \\
 & Token bottleneck, M = 512 & 6 & 98\%--100\% & 39\%--99\% & 52\%--98\% & 1\%--12\% & 85\%--95\% \\
 &  & 7 & 100\% & 80\%--100\% & 88\%--100\% & 2\%--19\% & 99.6\%--100\% \\
 &  & 8 & 100\% & 99.6\%--100\% & 99\%--100\% & 3\%--25\% & 100\% \\
 &  & 9 & 100\% & 100\% & 100\% & 4\%--32\% & 100\% \\
 & Token bottleneck, M = 1,024 & 6 & 99\%--100\% & 49\%--78\% & 58\%--77\% & 1\%--2\% & 80\%--89\% \\
 &  & 7 & 100\% & 86\%--99.6\% & 90\%--98\% & 1\%--3\% & 99\%--100\% \\
 &  & 8 & 100\% & 99\%--100\% & 98\%--100\% & 1\%--5\% & 100\% \\
 &  & 9 & 100\% & 100\% & 100\% & 2\%--8\% & 100\% \\
 & Nearest-neighbour field & 6 & 100\% & 100\% & 100\% & 90\% & 93\% \\
 &  & 7 & 100\% & 100\% & 100\% & 93\% & 96\% \\
 &  & 8 & 100\% & 100\% & 100\% & 96\% & 98\% \\
 &  & 9 & 100\% & 100\% & 100\% & 97\% & 98\% \\
 & Scaled polynomial & 6 & 100\% & 100\% & 100\% & 100\% & 100\% \\
 &  & 7 & 100\% & 100\% & 100\% & 100\% & 100\% \\
 &  & 8 & 100\% & 100\% & 100\% & 100\% & 100\% \\
 &  & 9 & 100\% & 100\% & 100\% & 100\% & 100\% \\
 & Zero field & 6 & 100\% & 100\% & 100\% & 100\% & 100\% \\
 &  & 7 & 100\% & 100\% & 100\% & 100\% & 100\% \\
 &  & 8 & 100\% & 100\% & 100\% & 100\% & 100\% \\
 &  & 9 & 100\% & 100\% & 100\% & 100\% & 100\% \\
\midrule
3D, finer mesh & Label-free, zero-shot & 6 & 100\% & 99.6\%--100\% & 100\% & 8\%--19\% & 92\%--96\% \\
 &  & 7 & 100\% & 100\% & 100\% & 9\%--20\% & 99\%--100\% \\
 &  & 8 & 100\% & 100\% & 100\% & 12\%--20\% & 100\% \\
 &  & 9 & 100\% & 100\% & 100\% & 14\%--21\% & 100\% \\
 & Supervised, displacement, zero-shot & 6 & 100\% & 100\% & 100\% & 82\%--97\% & 100\% \\
 &  & 7 & 100\% & 100\% & 100\% & 87\%--98\% & 100\% \\
 &  & 8 & 100\% & 100\% & 100\% & 91\%--99\% & 100\% \\
 &  & 9 & 100\% & 100\% & 100\% & 96\%--99\% & 100\% \\
 & Graph network, displacement, zero-shot & 6 & 100\% & 100\% & 100\% & 100\% & 100\% \\
 &  & 7 & 100\% & 100\% & 100\% & 100\% & 100\% \\
 &  & 8 & 100\% & 100\% & 100\% & 100\% & 100\% \\
 &  & 9 & 100\% & 100\% & 100\% & 100\% & 100\% \\
 & Nearest-neighbour field & 6 & 100\% & 100\% & 100\% & 100\% & 100\% \\
 &  & 7 & 100\% & 100\% & 100\% & 100\% & 100\% \\
 &  & 8 & 100\% & 100\% & 100\% & 100\% & 100\% \\
 &  & 9 & 100\% & 100\% & 100\% & 100\% & 100\% \\
 & Scaled polynomial & 6 & 100\% & 100\% & 100\% & 100\% & 100\% \\
 &  & 7 & 100\% & 100\% & 100\% & 100\% & 100\% \\
 &  & 8 & 100\% & 100\% & 100\% & 100\% & 100\% \\
 &  & 9 & 100\% & 100\% & 100\% & 100\% & 100\% \\
 & Label-free, then fine-mesh labels, 16 labels & 6 & 95\%--96\% & 100\% & 99.6\%--100\% & 16\%--43\% & 95\%--96\% \\
 &  & 7 & 99\%--100\% & 100\% & 100\% & 23\%--53\% & 99\%--100\% \\
 &  & 8 & 100\% & 100\% & 100\% & 26\%--63\% & 100\% \\
 &  & 9 & 100\% & 100\% & 100\% & 34\%--73\% & 100\% \\
 & Supervised from scratch, fine-mesh labels, 16 labels & 6 & 100\% & 100\% & 100\% & 100\% & 100\% \\
 &  & 7 & 100\% & 100\% & 100\% & 100\% & 100\% \\
 &  & 8 & 100\% & 100\% & 100\% & 100\% & 100\% \\
 &  & 9 & 100\% & 100\% & 100\% & 100\% & 100\% \\
 & Label-free, then fine-mesh labels, 64 labels & 6 & 66\%--72\% & 100\% & 80\%--99.6\% & 2\%--17\% & 89\%--96\% \\
 &  & 7 & 90\%--95\% & 100\% & 100\% & 4\%--23\% & 98\%--100\% \\
 &  & 8 & 98\%--99.6\% & 100\% & 100\% & 5\%--28\% & 99\%--100\% \\
 &  & 9 & 100\% & 100\% & 100\% & 9\%--38\% & 100\% \\
 & Supervised from scratch, fine-mesh labels, 64 labels & 6 & 100\% & 100\% & 100\% & 100\% & 100\% \\
 &  & 7 & 100\% & 100\% & 100\% & 100\% & 100\% \\
 &  & 8 & 100\% & 100\% & 100\% & 100\% & 100\% \\
 &  & 9 & 100\% & 100\% & 100\% & 100\% & 100\% \\
 & Token bottleneck, M = 512, zero-shot & 6 & 100\% & 100\% & 100\% & 38\%--61\% & 99\%--99.6\% \\
 &  & 7 & 100\% & 100\% & 100\% & 46\%--74\% & 100\% \\
 &  & 8 & 100\% & 100\% & 100\% & 59\%--80\% & 100\% \\
 &  & 9 & 100\% & 100\% & 100\% & 70\%--87\% & 100\% \\
 & Token bottleneck, M = 1,024, zero-shot & 6 & 100\% & 100\% & 100\% & 40\%--64\% & 98\%--99\% \\
 &  & 7 & 100\% & 100\% & 100\% & 59\%--68\% & 100\% \\
 &  & 8 & 100\% & 100\% & 100\% & 75\%--83\% & 100\% \\
 &  & 9 & 100\% & 100\% & 100\% & 80\%--93\% & 100\% \\
\midrule
3D, holdout & Label-free & 6 & 67\%--81\% & 1\%--3\% & 1\%--4\% & 0\%--0.1\% & 56\%--79\% \\
 &  & 7 & 91\%--97\% & 2\%--21\% & 12\%--28\% & 0\%--0.1\% & 84\%--95\% \\
 &  & 8 & 98\%--99.9\% & 12\%--71\% & 32\%--69\% & 0\%--0.3\% & 96\%--99\% \\
 &  & 9 & 99.6\%--100\% & 34\%--99.9\% & 67\%--96\% & 0\%--1\% & 99.7\%--100\% \\
 & Supervised, displacement & 6 & 92\%--95\% & 100\% & 100\% & 26\%--32\% & 100\% \\
 &  & 7 & 99\%--99.6\% & 100\% & 100\% & 32\%--35\% & 100\% \\
 &  & 8 & 99.9\%--100\% & 100\% & 100\% & 36\%--40\% & 100\% \\
 &  & 9 & 100\% & 100\% & 100\% & 44\%--46\% & 100\% \\
 & Graph network, displacement & 6 & 16\%--20\% & 100\% & 100\% & 26\%--32\% & 99\%--99.9\% \\
 &  & 7 & 62\%--68\% & 100\% & 100\% & 29\%--36\% & 100\% \\
 &  & 8 & 92\%--95\% & 100\% & 100\% & 33\%--39\% & 100\% \\
 &  & 9 & 99.6\%--99.7\% & 100\% & 100\% & 36\%--45\% & 100\% \\
 & Supervised, energy norm & 6 & 45\%--58\% & 0.1\%--1\% & 1\%--2\% & 0\%--0.1\% & 27\%--28\% \\
 &  & 7 & 82\%--88\% & 1\%--2\% & 3\%--7\% & 0\%--0.3\% & 52\%--55\% \\
 &  & 8 & 94\%--97\% & 2\%--4\% & 11\%--17\% & 0\%--0.3\% & 63\%--67\% \\
 &  & 9 & 98\%--99\% & 6\%--12\% & 31\%--44\% & 0\%--0.3\% & 77\%--84\% \\
\midrule
2D, validation & Label-free & 6 & 80\%--84\% & 16\%--20\% & 47\%--53\% & 0\%--0.4\% & 95\%--96\% \\
 &  & 8 & 96\%--97\% & 61\%--66\% & 82\%--86\% & 2\%--3\% & 98\%--99\% \\
 & Supervised, displacement & 6 & 58\%--67\% & 100\% & 100\% & 0.4\%--4\% & 100\% \\
 &  & 8 & 76\%--78\% & 100\% & 100\% & 2\%--8\% & 100\% \\
 & Supervised, energy norm & 6 & 44\%--73\% & 0.4\%--15\% & 2\%--45\% & 0\% & 24\%--89\% \\
 &  & 8 & 64\%--96\% & 2\%--55\% & 7\%--74\% & 0\%--2\% & 33\%--98\% \\
 & Graph network, displacement & 6 & 57\%--61\% & 100\% & 100\% & 29\%--32\% & 99.6\%--100\% \\
 &  & 8 & 74\%--75\% & 100\% & 100\% & 36\%--41\% & 100\% \\
 & Nearest-neighbour field & 6 & 100\% & 100\% & 100\% & 84\% & 96\% \\
 &  & 8 & 100\% & 100\% & 100\% & 92\% & 96\% \\
 & Scaled polynomial & 6 & 100\% & 100\% & 100\% & 100\% & 100\% \\
 &  & 8 & 100\% & 100\% & 100\% & 100\% & 100\% \\
 & Zero field & 6 & 100\% & 100\% & 100\% & 100\% & 100\% \\
 &  & 8 & 100\% & 100\% & 100\% & 100\% & 100\% \\
\midrule
2D, holdout & Label-free & 6 & 79\%--84\% & 17\%--19\% & 49\%--53\% & 0.1\%--0.2\% & 94\%--96\% \\
 &  & 8 & 94\%--95\% & 62\%--66\% & 81\%--85\% & 4\% & 99\% \\
 & Supervised, displacement & 6 & 56\%--63\% & 100\% & 100\% & 1\%--3\% & 99.98\%--100\% \\
 &  & 8 & 73\%--77\% & 100\% & 100\% & 4\%--8\% & 99.98\%--100\% \\
 & Supervised, energy norm & 6 & 41\%--72\% & 0.4\%--16\% & 2\%--46\% & 0.1\%--0.4\% & 21\%--87\% \\
 &  & 8 & 63\%--94\% & 2\%--59\% & 8\%--73\% & 0.2\%--4\% & 33\%--97\% \\
 & Graph network, displacement & 6 & 55\%--57\% & 100\% & 100\% & 29\%--33\% & 99.9\%--100\% \\
 &  & 8 & 71\%--72\% & 100\% & 100\% & 37\%--42\% & 99.98\%--100\% \\
\end{longtable}
\endgroup
\vspace{-\smallskipamount}
\noindent\begin{minipage}{\linewidth}\footnotesize\raggedright
The share of the instances on which the Ritz projection was strictly the more accurate than the model, for each seed of a trained network (the lowest and the highest over the three seeds; one value when they agree), and for the deterministic baselines. Higher is better in critical-region recall, lower in every other metric; ties, which occur only in critical-region recall, are counted with neither. A share strictly between 0\% and 100\% is never shown as either. Networks at 1,024 labelled instances; on the finer mesh the networks trained in-band are applied zero-shot and the few-shot networks are those of the transfer comparison. The shares for every other network and budget of the summaries are in the records.
\end{minipage}
\par\medskip

\end{landscape}

\section*{S2. The three-dimensional stiffness-norm run}

Table~\ref{tab:kmedians} gives every metric of the three transformers of the three-dimensional
stiffness-norm run on the 256 validation instances, as the seeds' medians over the
instances, the statistic that the run's pre-registered verdicts read, and as the seeds'
means. Table~\ref{tab:kcomparisons} gives the \numKdSecondary{} secondary comparisons that its
pre-registration fixed as exploratory readings, without a criterion: for each of the five
metrics, the stiffness-norm transformer against the supervised transformer trained on the
relative displacement error and against the label-free transformer, and the label-free
transformer against the supervised one, on the medians and on the means. Each reading
applies the noise guard of the paper: rel is the relative difference of the two networks'
seed-averaged values, and a difference counts only beyond the larger of 10\% and twice its
standard error over the three seeds. Without a true difference, a comparison can cross the
guard by chance, as the pre-registration's simulations of the rule show, so several
crossings among
\numKdSecondary{} are expected by chance. For the critical-region recall higher is better,
and the reading says which network was the better.

% generated from the committed records; do not edit
\begin{table}[tbp]
\centering
\footnotesize
\setlength{\tabcolsep}{4pt}
\caption{Three-dimensional stiffness-norm run: every metric of the three transformers on the 256 validation instances, the seeds' medians and the seeds' means over the instances (mean $\pm$ sample standard deviation over the three seeds).}
\label{tab:kmedians}
\resizebox{\ifdim\width>\textwidth\textwidth\else\width\fi}{!}{%
\begin{tabular}{llccccc}
\toprule
Network & Statistic & \shortstack{Displacement \\ error} & \shortstack{Relative \\ energy gap} & \shortstack{von Mises \\ error} & \shortstack{Peak stress \\ error} & \shortstack{Critical-region \\ recall} \\
\midrule
Label-free & seed medians & 0.0271 $\pm$ 0.0022 & 0.00752 $\pm$ 0.00176 & 0.0527 $\pm$ 0.0040 & 0.0785 $\pm$ 0.0149 & 0.790 $\pm$ 0.023 \\
 & seed means & 0.0300 $\pm$ 0.0028 & 0.00912 $\pm$ 0.00228 & 0.0554 $\pm$ 0.0045 & 0.0807 $\pm$ 0.0134 & 0.793 $\pm$ 0.022 \\
\midrule
Supervised ($\mathcal{L}_D$) & seed medians & 0.0415 $\pm$ 0.0026 & 0.251 $\pm$ 0.026 & 0.259 $\pm$ 0.015 & 0.222 $\pm$ 0.014 & 0.628 $\pm$ 0.009 \\
 & seed means & 0.0470 $\pm$ 0.0038 & 31.2 $\pm$ 53.5 & 0.370 $\pm$ 0.140 & 0.784 $\pm$ 0.807 & 0.619 $\pm$ 0.005 \\
\midrule
Supervised, stiffness norm & seed medians & 0.0203 $\pm$ 0.0019 & 0.00444 $\pm$ 0.00010 & 0.0425 $\pm$ 0.0013 & 0.0379 $\pm$ 0.0008 & 0.858 $\pm$ 0.001 \\
 & seed means & 0.0227 $\pm$ 0.0022 & 0.00504 $\pm$ 0.00022 & 0.0440 $\pm$ 0.0014 & 0.0386 $\pm$ 0.0011 & 0.858 $\pm$ 0.001 \\
\bottomrule
\end{tabular}}
\end{table}


% generated from the committed records; do not edit
\begin{table}[tbp]
\centering
\footnotesize
\setlength{\tabcolsep}{4pt}
\caption{Three-dimensional stiffness-norm run: the 30 secondary comparisons, exploratory, each with the noise guard, on the seeds' medians and on the seeds' means.}
\label{tab:kcomparisons}
\resizebox{\ifdim\width>\textwidth\textwidth\else\width\fi}{!}{%
\begin{tabular}{llcccc}
\toprule
Comparison (A against B) & Metric & Medians: rel (guard) & Reading & Seed means: rel (guard) & Reading \\
\midrule
Stiffness norm against supervised ($\mathcal{L}_D$) & Displacement error & $-$51.1\% (guard 10.0\%) & A better & $-$51.8\% (guard 10.8\%) & A better \\
 & Relative energy gap & $-$98.2\% (guard 11.7\%) & A better & $-$99.98\% (guard 197.9\%) & within \\
 & von Mises error & $-$83.6\% (guard 10.0\%) & A better & $-$88.1\% (guard 43.7\%) & A better \\
 & Peak stress error & $-$82.9\% (guard 10.0\%) & A better & $-$95.1\% (guard 118.9\%) & within \\
 & Critical-region recall & +36.7\% (guard 10.0\%) & A better & +38.6\% (guard 10.0\%) & A better \\
\midrule
Stiffness norm against label-free & Displacement error & $-$25.0\% (guard 12.3\%) & A better & $-$24.3\% (guard 13.7\%) & A better \\
 & Relative energy gap & $-$41.0\% (guard 27.0\%) & A better & $-$44.7\% (guard 29.0\%) & A better \\
 & von Mises error & $-$19.4\% (guard 10.0\%) & A better & $-$20.6\% (guard 10.0\%) & A better \\
 & Peak stress error & $-$51.7\% (guard 22.0\%) & A better & $-$52.1\% (guard 19.2\%) & A better \\
 & Critical-region recall & +8.7\% (guard 10.0\%) & within & +8.1\% (guard 10.0\%) & within \\
\midrule
Label-free against supervised ($\mathcal{L}_D$) & Displacement error & $-$34.7\% (guard 10.0\%) & A better & $-$36.3\% (guard 11.6\%) & A better \\
 & Relative energy gap & $-$97.0\% (guard 11.8\%) & A better & $-$99.97\% (guard 197.9\%) & within \\
 & von Mises error & $-$79.6\% (guard 10.0\%) & A better & $-$85.0\% (guard 43.7\%) & A better \\
 & Peak stress error & $-$64.6\% (guard 10.6\%) & A better & $-$89.7\% (guard 118.9\%) & within \\
 & Critical-region recall & +25.8\% (guard 10.0\%) & A better & +28.2\% (guard 10.0\%) & A better \\
\bottomrule
\end{tabular}}
\end{table}


\section*{S3. Timing: threads and an earlier measurement}

The solvers of the cost table ran with the machine's \numCostThreads{} threads. The effect of
the number of threads was measured on another host of the same type, with
\numThreadsHost{} threads, on the first \numThreadsValN{} validation and \numThreadsFineN{}
fine-mesh instances of the three-dimensional comparison: with one thread, the direct solve
took the same time, to within \numThreadsDirectDev{}, and conjugate gradients
\numThreadsCgGainMin{} to \numThreadsCgGainMax{} less; a repeat with all threads agreed with
the first measurement to within \numThreadsDrift{}. In two dimensions (\numThreadsTwoN{}
instances, without a repeat), one thread was \numThreadsTwoDirect{} faster for the direct
solve and \numThreadsTwoCg{} for conjugate gradients.


The timings of the cost table of the paper were measured on 28 October 2026, on the machine
that computed the Ritz projection, so that every row of the table comes from one machine.
The same instances and networks had been timed on 8 October 2026 on another machine
(\numRerunCpu{}, \numRerunThreads{} threads). The surrogate's accuracies agreed between the
two measurements to within a relative \numRerunAccDev{}, and the iteration counts of
conjugate gradients differed by at most \numRerunIters{} per load case. The times of the
earlier measurement are committed with the records.

\end{document}
